%%%%%%%%%%  scan idx 291  |  folio 278  %%%%%%%%%%
on
peut écrire, comme on le voit facilement, \(x=\cl(E')-\cl(E'')\) pour \(E',E''\in\Ob C\) (on peut même s'arranger, dans le cas non respé, pour que \(E''=0\)). Comme la filtration \((C_i)\) est exhaustive, il existe un entier \(i\geq0\) tel que \(E'\) et \(E''\) appartiennent à \(C_i\); on a alors \(x\in k(C)_i=k(C)_0\) et l'on gagne.

\textbf{Exercice 1.9.} Soient \(C\) une catégorie abélienne, \(A\) une sous-catégorie pleine stable par sommes finies, sous-objets et quotients, et telle que \(\Ob A\) soit dense dans \(\Ob C\). Alors l'homomorphisme \(k(A,C)\longrightarrow k(C)\) défini par la fonction additive \(E\mapsto\cl_C(E)\) sur \(\Ob A\) (cf. (1.5)) est un isomorphisme.

\emph{Hint} : Poser \(A=A_0\), et définir \(A_i\), pour \(i\geq1\), comme la sous-catégorie pleine de \(C\) formée des objets de \(C\) qui sont extensions de deux objets de \(C_{i-1}\). Prouver que \(A_i\) vérifie les mêmes hypothèses que \(A\), puis montrer que l'homomorphisme canonique \(\varinjlim k(A_i,C)\longrightarrow k(C)\) est un isomorphisme (la conclusion voulue en résulte aussitôt).

\section*{2. Les foncteurs \(K_{\bullet}\) et \(K^{\bullet}\) d'un topos annelé}

\textbf{2.1.} Soit \((X,A)\) un topos annelé. Rappelons les notations suivantes :

\(D^b({}_A X)_{\coh}\) = catégorie des complexes pseudo-cohérents et à cohomologie bornée de \(A\)-Modules à gauche (I 2.15).

\(D({}_A X)_{\uparf}\) = catégorie des complexes de \(A\)-Modules à gauche d'amplitude parfaite finie (I 5.8.4).

\(\Loclib({}_A X)\) = catégorie des \(A\)-Modules à gauche localement facteurs directs de Modules libres de type fini (II 2.0).

\(\Coh({}_A X)\) = catégorie des \(A\)-Modules à gauche cohérents (relativement à la catégorie des \(A\)-Modules libres de type fini) (I 3.1).

